\documentclass[10pt,a4paper]{amsart}
\usepackage[margin=19mm]{geometry}
\usepackage{amsmath,amssymb,mathtools}
\usepackage[scr=rsfs,cal=euler]{mathalfa}
\usepackage{xfrac}
\usepackage[unicode,hidelinks,bookmarks=false]{hyperref}
\newcommand{\ie}{i.e.\ }
\newcommand{\cf}{cf.\ }
\newcommand{\resp}{respectively\ }
\newcommand{\Subsection}{\subsection}
\providecommand{\nobreakdash}{\nobreak\hbox{-}}
\setlength{\emergencystretch}{2em}
\multlinegap0pt
\usepackage{bm}
\usepackage{bbm}
\usepackage{dsfont}
\usepackage{xspace}
\usepackage{cancel}
\usepackage{mathtools}
\usepackage{amsthm}
\makeatletter
\@ifundefined{theorem}{\newtheorem{theorem}{Theorem}}{}
\@ifundefined{lemma}{\newtheorem{lemma}[theorem]{Lemma}}{}
\@ifundefined{proposition}{\newtheorem{proposition}[theorem]{Proposition}}{}
\makeatother
\newcommand{\Z}{\mathbb{Z}}
\newcommand{\disc}{\text{disc}}
\title[Certain Siegel Cusp Forms with Level are Determined by their]{Certain Siegel Cusp Forms with Level are Determined by their Fundamental Fourier Coefficients}
\author{Sidney Washburn}
\date{Source: arXiv:2507.17002v2 (CC BY 4.0)}
\begin{document}
\maketitle

\noindent\emph{Source and licence.} Certain Siegel Cusp Forms with Level are Determined by their Fundamental Fourier Coefficients. Version arXiv:2507.17002v2 (CC BY 4.0), distributed under the Creative Commons Attribution 4.0 International licence, as recorded in the arXiv licence metadata for this exact version and shown on the abstract page https://arxiv.org/abs/2507.17002v2. Full rights notice: licenses/arxiv-2507.17002-CC-BY-4.0.txt.\par\smallskip

\noindent\textit{Editorial scope.} Counted unit: the Proposition whose source label is \texttt{nonvanishing of twist} in the article (source0.tex lines 689--691, source label \texttt{nonvanishing of twist}), delivered together with the article's own complete original proof of it (source0.tex lines 693--714, one \texttt{proof} environment of 22 lines). No mathematics is written, repaired or re-derived: statement and proof are the article's own text line for line. Required context retained uncounted: 3.3 (lines 365--367); primitivity lemma (lines 557--559). Every definition, display and label of those lines is retained unchanged. Omitted from this package and not redistributed: 1--364, 368--556, 560--688, 692--692, 715--953. Nothing inside a retained excerpt is omitted or altered: every retained body is delivered exactly as the article's lines are written, so no omission receipt of any kind accompanies this package. Citations: the retained excerpts carry no citation command at all, so no bibliography entry of the article is transcribed and no reference is redistributed. Reference adaptation: none was needed and none was made; every reference inside the delivered unit resolves inside it, so no unresolved reference or citation is produced. Printed slips kept literal: no printed word, symbol, hypothesis or number of the counted statement or of its proof was altered. Layout and class: the article's own class is \texttt{article} with packages including calc, amssymb, amsmath, ulem, stmaryrd, fullpage, amsthm, enumerate, alltt, color, graphicx, xy, verbatim, geometry; the wrapper is the lane's pinned \texttt{amsart} editorial wrapper at 10pt on A4, which restates the theorem-like environments the retained text needs and adds the article's own macro definitions that the retained text needs (\texttt{\textbackslash Z}, \texttt{\textbackslash disc}), with extra wrapper packages (bm, bbm, dsfont, xspace, cancel, mathtools). The delivered numbering is the unit's own. Assets: no retained span contains a figure environment, a graphics-inclusion command, a PDF-page inclusion, a table or a TikZ picture; no image file is copied into this package, the package declares no asset dependency and no image file is redistributed with it. Licence: version arXiv:2507.17002v2 of this article is distributed under the Creative Commons Attribution 4.0 International licence, as recorded in the arXiv licence metadata for this exact version and shown on the abstract page https://arxiv.org/abs/2507.17002v2. A full rights notice is delivered at licenses/arxiv-2507.17002-CC-BY-4.0.txt. Characters: every retained excerpt is pure ASCII, so the delivered engine, XeLaTeX with its default Latin Modern text font, reports no missing character.
\begin{theorem}\label{3.3}
    Let $\varphi \in J^{\text{cusp}}_{k, T}(\Gamma_0(N),\chi)$ be a nonzero Jacobi form with fundamental index $T \in \Lambda_n^+$, nebentypus $\chi$, and odd level $N$. If $(\disc(T),N) = 1$, then there exists $\mu \in \Z^n/2T\Z^n$ with so that $h_\mu \neq 0$ and $T^{-1}[\mu/2]$ has maximum possible denominator, equal to $\mathrm{lvl}(T)$.
\end{theorem} 

\begin{lemma} \label{primitivity lemma}
    Let $T \in \Lambda_n^+$ be fundamental with discriminant $d_T$. Fix any $\Z$-module isomorphism $\alpha:\Z^n/(2T)\Z^n \to \Z/d'\Z$, where $d' = d_T$ if $n$ is even and $d' = 2d_T$ if $n$ is odd. Then $\mu \in \Z^n/(2T)\Z^n$ is primitive if and only if $\alpha(\mu) \in (\Z/d'\Z)^\times$. 
\end{lemma}

\begin{proposition} \label{nonvanishing of twist}
    Let $\phi \in J_{k,T}(\Gamma_0(N),\chi)$ be a nonzero Jaccobi form of fundamental index $T$. Assume that $(d_T,N) = 1$ and that $N$ is odd and square-free. Then there exists a Dirichlet character $\epsilon$ modulo $d'$ so that the corresponding twisted Eichler--Zagier map $h_\epsilon$ does not vanish. Here, $d' = d_T$ or $2d_T$ just as in Lemma \ref{primitivity lemma}.
\end{proposition}

\begin{proof}
    Suppose $h_\epsilon = 0$ for every $\epsilon$. This is equivalent to the vanishing of the following expression:

    \[
    \left( \epsilon(\mu) \right)_{\epsilon,\mu} \cdot (h_\mu)_\mu = A(d') \cdot (h_\mu)_\mu
    \]

    Here, the matrix $A(d')$ has rows indexed by Dirichlet characters $\epsilon$ modulo $d'$ and columns indexed by primitive components $\mu \in \Z^n/(2T)\Z^n$. By Lemma \ref{primitivity lemma}, there are $\phi(d')$ columns, which is equal to the number of rows. So $A(d')$ is a square matrix. From Theorem $\ref{3.3}$, the column $(h_\mu)_\mu$ does not vanish. Hence, we will have reached a contradiction if $A(d')$ is invertible. We prove this statement now. \\

    We induct on the number of prime factors of $d'$. When $d' = 1$ or $2$, $A(d')$ is simply the number $1$ and this is obvious. Suppose that $d' = p$ is an odd prime. Then, up to rearranging rows and columns, $A(p)$ takes the following form: 

    \[
    A(p) = \begin{pmatrix}
        1 & 1 & 1 & \cdots & 1 \\
        1 & \zeta_{p-1} & \zeta_{p-1}^2 & \cdots & \zeta_{p-1}^{p-2} \\ 1 & \zeta_{p-1}^2 & (\zeta_{p-1}^2)^2 & \cdots & (\zeta_{p-1}^2)^{p-2} \\ \vdots & \vdots & \ddots & \cdots & \vdots\\ 1 & \zeta_{p-1}^{p-2} & (\zeta_{p-1}^{p-2})^2 & \cdots & (\zeta_{p-1}^{p-2})^{p-2}
    \end{pmatrix}
    \]

    This follows by noting that $(\Z/p\Z)^\times$ is cyclic of order $p-1$. So $A(p)$ is a Vandermonde matrix with determinant $\displaystyle \prod_{0 \leq i < j \leq p-2} (\zeta_{p-1}^i - \zeta_{p-1}^j)$. Of course, this is nonzero and so $A(p)$ is invertible. \\

    In general, let $d' = p_1 \cdots p_t$ be the prime factorization of $d'$. By the Chinese Remainder Theorem, we have $A(d') = A(d'/p_t) \otimes A(p_t)$. By induction on $t$ along with the well-known rank identity $\text{rank}(A \otimes B) = \text{rank}(A) \cdot \text{rank}(B)$, $A(d')$ is invertible and the proof is complete. 
\end{proof}

\end{document}
